\subsection{Wellencharakter des Elektrons (de-Broglie-Hypothese)}
\label{subsec:wellencharakter_elektron}

Die klassische Physik beschreibt das Elektron\index{Elektron} zunächst als
kleines geladenes Teilchen\index{Teilchen}. Es besitzt eine Masse, eine
elektrische Ladung und einen Impuls\index{Impuls}. Mit diesen Größen lässt sich
seine Bewegung in elektrischen und magnetischen Feldern gut beschreiben. Diese
Vorstellung war jedoch nur ein Teil der Wirklichkeit.

Zu Beginn des 20. Jahrhunderts wurde deutlich, dass Licht nicht nur
Welleneigenschaften\index{Welleneigenschaft} besitzt, sondern auch
Teilcheneigenschaften\index{Teilcheneigenschaft}. Beim Photoeffekt
\index{Photoeffekt} zeigte sich, dass Licht seine Energie in einzelnen Quanten
abgibt. Diese Lichtquanten werden Photonen\index{Photon} genannt. Damit entstand
eine neue Frage: Wenn Licht sowohl Wellen- als auch Teilcheneigenschaften
besitzt, könnte dann auch Materie Welleneigenschaften besitzen?

Louis de Broglie\index{de Broglie, Louis} formulierte 1924 eine weitreichende
Hypothese. Er ordnete jedem Teilchen mit einem Impuls eine Wellenlänge zu. Für
ein Teilchen mit dem Impuls \(p\) gilt:

\begin{equation}
	\lambda = \frac{h}{p}
	\label{eq:de_broglie_wellenlaenge}
\end{equation}

Dabei ist \(\lambda\)\index{Wellenlänge} die de-Broglie-Wellenlänge
\index{de-Broglie-Wellenlänge}, \(h\)\index{Plancksches Wirkungsquantum} das
Plancksche Wirkungsquantum und \(p\) der Impuls des Teilchens. Für ein
nichtrelativistisches Elektron gilt näherungsweise:

\begin{equation}
	p = m_\mathrm{e} v
	\label{eq:impuls_elektron_klassisch}
\end{equation}

Damit ergibt sich für die Wellenlänge eines Elektrons:

\begin{equation}
	\lambda = \frac{h}{m_\mathrm{e} v}
	\label{eq:de_broglie_elektron}
\end{equation}

Diese Gleichung zeigt einen entscheidenden Zusammenhang: Je größer der Impuls
eines Elektrons ist, desto kleiner ist seine Wellenlänge. Ein langsam bewegtes
Elektron besitzt daher eine größere Wellenlänge als ein schnell bewegtes
Elektron.

Für Elektronen in atomaren Größenordnungen ist diese Wellenlänge nicht
vernachlässigbar. Sie liegt in einem Bereich, in dem Beugung\index{Beugung}
und Interferenz\index{Interferenz} auftreten können. Genau deshalb kann ein
Elektron nicht mehr vollständig als klassisches Teilchen mit einer exakt
bestimmten Bahn beschrieben werden.

\begin{PhysicsBox}[Bedeutung der de-Broglie-Hypothese]
	\label{box:broglie_hypothese}
	Die de-Broglie-Hypothese verbindet den klassischen Impuls eines Teilchens mit
	einer Wellenlänge. Damit wird Materie nicht durch eine klassische Bahn allein
	beschrieben, sondern erhält einen wellenartigen Charakter.
\end{PhysicsBox}

Die Welle, die einem Elektron zugeordnet wird, ist jedoch keine gewöhnliche
mechanische Welle wie eine Wasserwelle oder Schallwelle. Sie beschreibt nicht
eine Auslenkung eines materiellen Mediums. In der Quantenmechanik führt dieser
Gedanke zur Wellenfunktion\index{Wellenfunktion}. Aus ihr lassen sich
Wahrscheinlichkeiten\index{Wahrscheinlichkeit} für mögliche Messergebnisse
ableiten.

Der Wellencharakter des Elektrons\index{Wellencharakter!des Elektrons} ist
daher kein anschauliches Zusatzbild, sondern eine grundlegende Eigenschaft der
quantenmechanischen Beschreibung. Er erklärt, warum Elektronen Beugungs- und
Interferenzerscheinungen zeigen können und warum die klassische Vorstellung
fester Elektronenbahnen im Atom nicht ausreicht.